\chapter{Abdominal Pain}

\section*{Definition}
Abdominal pain is discomfort or pain anywhere between the lower chest and groin. It may arise from the abdominal organs, abdominal wall, pelvis, urinary tract, or structures outside the abdomen such as the heart or lungs. Its location, character, severity, timing, and associated symptoms can help identify the cause, but pain intensity alone does not reliably distinguish serious from minor illness.

\section*{When to See a Doctor}
\begin{itemize}
 \item Seek emergency care for sudden or severe pain, a rigid or markedly tender abdomen, fainting, confusion, shock-like symptoms, persistent vomiting or inability to keep fluids down, vomiting blood, black or bloody stools, chest pain, trouble breathing, inability to pass stool or gas especially with vomiting, or possible pregnancy with pain or bleeding.
 \item Arrange prompt medical assessment for fever or feeling very unwell, jaundice, dehydration, persistent or recurrent pain, pain that is worsening or localized to one area, recent abdominal injury or surgery, or severe diarrhea.
\end{itemize}

\section*{How It Develops}
Location, character, timing, and triggers provide diagnostic clues. Visceral pain arises from internal organs and is often vague or poorly localized; parietal or somatic pain from the abdominal lining or wall is usually sharper and more localized. Referred pain is felt away from its source, such as shoulder pain from irritation of the diaphragm. These patterns are clues rather than a diagnosis, and serious illness can present atypically.

\section*{Causes}
\begin{itemize}
 \item Gastroenteritis, food poisoning, constipation, gas, or functional gastrointestinal disorders such as IBS and functional dyspepsia
 \item Appendicitis, diverticulitis, inflammatory bowel disease, intestinal obstruction, or bowel ischemia
 \item Peptic ulcer disease, reflux, pancreatitis, or gastrointestinal bleeding
 \item Gallstones, cholecystitis, hepatitis, or other liver disease
 \item Urinary tract infection, kidney stone, or kidney infection
 \item In people with reproductive organs: ectopic pregnancy, miscarriage, ovarian cyst or torsion, pelvic inflammatory disease, endometriosis, or labor-related conditions
 \item Hernia, abdominal-wall strain, shingles, or abdominal injury
 \item Aortic aneurysm, heart attack, pneumonia, pulmonary embolism, or metabolic conditions such as diabetic ketoacidosis can sometimes cause abdominal or upper-abdominal pain
\end{itemize}

\section*{Related Symptoms}
\begin{itemize}
 \item Nausea
 \item Vomiting
 \item Bloating
 \item Loss of appetite
 \item Constipation or diarrhea
 \item Fever, chills, jaundice, urinary symptoms, vaginal bleeding, or abnormal discharge when relevant
\end{itemize}

\section*{Medical Evaluation}
\begin{itemize}
 \item A focused history assesses onset, location, quality, severity, duration, radiation, aggravating and relieving factors, relation to meals or bowel movements, associated symptoms, medication and alcohol use, previous operations, medical history, and pregnancy possibility when relevant.
 \item Physical examination includes vital signs and assessment of hydration, abdominal tenderness, guarding, rigidity, distension, bowel sounds, hernias, and masses. A pelvic, rectal, chest, or cardiac examination may be appropriate based on the symptoms.
 \item Tests may include a complete blood count, electrolytes, kidney function, liver tests, lipase, inflammatory markers, urinalysis and urine culture, stool tests, or pregnancy testing when indicated. An ECG and cardiac blood tests may be needed for upper-abdominal or epigastric pain with cardiac risk or associated chest symptoms.
 \item Imaging is selected according to the suspected cause, age, pregnancy status, and examination. Ultrasound is useful for biliary, pelvic, urinary, and some vascular conditions. CT of the abdomen and pelvis is often appropriate for acute nonlocalized pain or suspected obstruction, inflammation, or perforation; MRI may be preferred in selected patients, including some pregnant patients. Necessary imaging should not be withheld solely because pregnancy is possible.
\end{itemize}

\section*{Treatment Options}
\begin{itemize}
 \item Treatment depends on the cause and may include oral or intravenous fluids, appropriate pain relief, antiemetics, acid suppression, antibiotics for selected infections, treatment of constipation, or urgent surgery or other procedures.
 \item Suspected appendicitis, obstruction, perforation, ectopic pregnancy, ischemia, or internal bleeding requires urgent clinician-directed treatment.
\end{itemize}

\section*{Self-Care and Home Management}
\begin{itemize}
 \item For mild, improving pain without warning signs, drink fluids and eat small meals as tolerated. Oral rehydration may help when mild vomiting or diarrhea is present.
 \item Avoid alcohol and foods that clearly worsen symptoms.
 \item Acetaminophen may help when used exactly as labelled, but ask a clinician or pharmacist first if you have liver disease, drink heavily, or take other products containing acetaminophen. Avoid NSAIDs if ulcer, kidney disease, bleeding risk, or pregnancy is possible unless a clinician advises otherwise.
 \item Do not use laxatives or anti-diarrheal medicines when obstruction or an inflammatory/infectious cause is possible without medical advice.
\end{itemize}

\section*{References}
\begin{thebibliography}{99}
\bibitem{abdominal_pain:ref1} Cartwright SL, Knudson MP. Evaluation of acute abdominal pain in adults. \textit{Am Fam Physician}. 2015;92(10):902--908.
\bibitem{abdominal_pain:ref2} American College of Radiology. ACR Appropriateness Criteria: Acute Nonlocalized Abdominal Pain. Current online guidance.
\bibitem{abdominal_pain:ref3} National Institute of Diabetes and Digestive and Kidney Diseases. Diagnosis of Appendicitis. Reviewed 2021.
\bibitem{abdominal_pain:ref4} MedlinePlus. Abdominal Pain. U.S. National Library of Medicine. Accessed 2026.
\bibitem{abdominal_pain:ref5} American College of Obstetricians and Gynecologists. Diagnostic Imaging During Pregnancy and Lactation. ACOG Committee Opinion No. 723, reaffirmed 2021.
\end{thebibliography}
